\chapter{svnrdump Reference---Remote Subversion Repository Data
Migration}\label{svn.ref.svnrdump}

\texttt{svnrdump} joined the Subversion tool chain in the Subversion 1.7
release. It is best described as a network-aware version of the
\texttt{svnadmin\ dump} and \texttt{svnadmin\ load} commands, paired
together and released as a separate standalone program. We discuss the
process of dumping and loading repository data---using both
\texttt{svnadmin} and \texttt{svnrdump}--- in
\hyperref[svn.reposadmin.maint.migrate]{Migrating Repository Data
Elsewhere}.

Options in \texttt{svnrdump} are global, just as they are in
\texttt{svn} and \texttt{svnadmin}:

\protect\phantomsection\label{svn.ref.svnrdump.sw}{}

\begin{description}
\item[\protect\phantomsection\label{svn.ref.svnrdump.sw.config_dir}{}\texttt{-\/-config-dir}
\textless DIR\textgreater{}]
Instructs Subversion to read configuration information from the
specified directory instead of the default location
(\texttt{.subversion} in the user\textquotesingle s home directory).
\item[\protect\phantomsection\label{svn.ref.svnrdump.sw.config_option}{}\texttt{-\/-config-option}
\textless FILE\textgreater:\textless SECTION\textgreater:\textless OPTION\textgreater={[}\textless VALUE\textgreater{]}]
Sets, for the duration of the command, the value of a runtime
configuration option. \textless FILE\textgreater{} and
\textless SECTION\textgreater{} are the runtime configuration file
(either \texttt{config} or \texttt{servers}) and the section thereof,
respectively, which contain the option whose value you wish to change.
\textless OPTION\textgreater{} is, of course, the option itself, and
\textless VALUE\textgreater{} the value (if any) you wish to assign to
the option. For example, to temporarily disable the use of the
compression in the HTTP protocol, use
\texttt{-\/-config-option=servers:global:http-compression=no}. You can
use this option multiple times to change multiple option values
simultaneously.
\item[\protect\phantomsection\label{svn.ref.svnrdump.sw.incremental}{}\texttt{-\/-incremental}]
Dump a revision or revision range only as a diff against the previous
revision, instead of the default, which is begin a dumped revision range
with a complete expansion of all contents of the repository as of that
revision.
\item[\protect\phantomsection\label{svn.ref.svnrdump.sw.no_auth_cache}{}\texttt{-\/-no-auth-cache}]
Prevents caching of authentication information (e.g., username and
password) in the Subversion runtime configuration directories.
\item[\protect\phantomsection\label{svn.ref.svnrdump.sw.non_interactive}{}\texttt{-\/-non-interactive}]
In the case of an authentication failure or insufficient credentials,
prevents prompting for credentials (e.g., username or password). This is
useful if you\textquotesingle re running Subversion inside an automated
script and it\textquotesingle s more appropriate to have Subversion fail
than to prompt for more information.
\item[\protect\phantomsection\label{svn.ref.svnrdump.sw.password}{}\texttt{-\/-password}
\textless PASSWD\textgreater{}]
Specifies the password to use when authenticating against a Subversion
server. If not provided, or if incorrect, Subversion will prompt you for
this information as needed.
\item[\protect\phantomsection\label{svn.ref.svnrdump.sw.quiet}{}\texttt{-\/-quiet}
(\texttt{-q})]
Requests that the client print only essential information while
performing an operation.
\item[\protect\phantomsection\label{svn.ref.svnrdump.sw.revision}{}\texttt{-\/-revision}
(\texttt{-r}) \textless ARG\textgreater{}]
Specifies a particular revision or revision range on which to operate.
\item[\protect\phantomsection\label{svn.ref.svnrdump.sw.trust_server_cert}{}\texttt{-\/-trust-server-cert}]
Used with \texttt{-\/-non-interactive} to accept any unknown SSL server
certificates without prompting.
\item[\protect\phantomsection\label{svn.ref.svnrdump.sw.username}{}\texttt{-\/-username}
\textless NAME\textgreater{}]
Specifies the username to use when authenticating against a Subversion
server. If not provided, or if incorrect, Subversion will prompt you for
this information as needed.
\end{description}

\section{svnrdump dump}\label{svn.ref.svnrdump.c.dump}

\emph{}

\textbf{Synopsis}

\texttt{svnrdump\ dump\ SOURCE\_URL}

\subsection{Description}

Dump---that is, generate a repository dump stream of---revisions of the
repository item located at \textless SOURCE\_URL\textgreater, printing
the information to standard output. By default, the entire history will
be included in the dump stream, but the scope of the operation can be
limited via the use of the \texttt{-\/-revision} (\texttt{-r}) option.

\subsection{Options}

\begin{verbatim}
--config-dir DIR
--config-option FILE:SECTION:OPTION=[VALUE]
--file (-F) FILE
--force-interactive
--incremental
--no-auth-cache
--non-interactive
--password PASSWD
--password-from-stdin
--quiet (-q)
--revision (-r) ARG
--trust-server-cert
--trust-server-cert-failures ARG
--username NAME
\end{verbatim}

\subsection{Examples}

Generate a dump stream of the full history of a remote repository
(assuming that the user as who this runs has authorization to read all
paths in the repository).

\begin{verbatim}
$ svnrdump dump http://svn.example.com/repos/calc > full.dump
* Dumped revision 0.
* Dumped revision 1.
* Dumped revision 2.
…
\end{verbatim}

Incrementally dump a single revision from that same repository:

\begin{verbatim}
$ svnrdump dump http://svn.example.com/repos/calc \
           -r 21 --incremental > inc.dump
* Dumped revision 21.
$
\end{verbatim}

\section{svnrdump help}\label{svn.ref.svnrdump.c.help}

\emph{Help!}

\textbf{Synopsis}

\texttt{svnrdump\ help}

\subsection{Description}

Use this to get help. Well, certain kinds of help. Need white lab coat
and straightjackets kind of help? There\textquotesingle s a whole
different protocol for that sort of thing.

\subsection{Options}

None

\section{svnrdump load}\label{svn.ref.svnrdump.c.load}

\emph{}

\textbf{Synopsis}

\texttt{svnrdump\ load\ DEST\_URL}

\subsection{Description}

Read from standard input revision information described in a Subversion
repository dump stream, and load that information into the repository
located at \textless DEST\_URL\textgreater.

\subsection{Options}

\begin{verbatim}
--config-dir DIR
--config-option FILE:SECTION:OPTION=[VALUE]
--file (-F) FILE
--force-interactive
--no-auth-cache
--non-interactive
--password PASSWD
--password-from-stdin
--quiet (-q)
--skip-revprop
--trust-server-cert
--trust-server-cert-failures ARG
--username NAME
\end{verbatim}

\subsection{Examples}

Dump the contents of a local repository, and use \texttt{svnrdump\ load}
to transfer that revision information into an existing remote
repository:

\begin{verbatim}
$ svnadmin dump -q /var/svn/repos/new-project | \
      svnrdump load http://svn.example.com/repos/new-project
* Loaded revision 0
* Loaded revision 1.
* Loaded revision 2.
…
\end{verbatim}

To operate properly \texttt{svnrdump\ load} requires that the target
repository have revision property modification enabled via the
pre-revprop-change hook. For details about that hook, see
\hyperref[svn.ref.reposhooks.pre-revprop-change]{pre-revprop-change}.

